\section{实验环境与问题建模}

本章把报告后续使用的量子信息概念和路由编译模型放在同一条线上说明。量子线路路由
不是普通图搜索的直接套用：算法每插入一个 SWAP，都对应量子态在物理载体上的交换；
算法每推进一个双比特门，都必须同时满足线路依赖关系和芯片耦合约束。把这一层关系
讲清楚，后面的 SABRE 复现、NPQR 设计和实验指标才有明确的物理含义。

\subsection{量子线路基础}

一个量子比特的纯态写作
\[
  |\psi\rangle=\alpha|0\rangle+\beta|1\rangle,\qquad
  |\alpha|^2+|\beta|^2=1 .
\]
这里 $|0\rangle$ 和 $|1\rangle$ 是计算基，$\alpha$、$\beta$ 是复振幅。测量得到
$0$ 或 $1$ 的概率分别为 $|\alpha|^2$ 和 $|\beta|^2$。多个量子比特组成寄存器后，
状态空间由张量积给出，$n$ 个量子比特对应 $2^n$ 维复向量空间。量子线路中的门操作
不是任意线性变换，而是保持内积和概率归一性的酉变换。

这一定义决定了路由编译的基本边界。编译器可以改变逻辑量子比特放置在哪个物理位置，
也可以插入语义明确的辅助门，但不能改变原线路描述的量子演化。换言之，路由输出的
线路必须和输入线路在逻辑语义上等价，只是执行时使用了不同的物理量子比特顺序和
额外的交换操作。

\subsubsection*{量子门与 SWAP}

单比特门作用在一个二维子空间上，常见例子包括 $X$、$H$ 和相位门。双比特门作用在
两个量子比特的张量积空间上，是产生和调控纠缠的关键。本项目关注的硬件约束主要
来自双比特门，因为真实芯片通常只允许相邻物理量子比特之间直接执行 CNOT 或 CZ。

SWAP 是路由编译中最重要的辅助门。它的作用可以写为
\[
  \mathrm{SWAP}|a\rangle|b\rangle=|b\rangle|a\rangle .
\]
在矩阵形式下，SWAP 是酉矩阵；在物理解释上，它交换两个物理位置上承载的逻辑量子
态。因此，若逻辑量子比特 $q_i$ 和 $q_j$ 当前不相邻，编译器可以沿硬件边插入若干
SWAP，把二者移动到相邻位置后再执行目标双比特门。这个过程保持理想线路语义，却
增加门数和深度。

\subsection{线路与硬件拓扑}

量子线路可以看成有序门列
\[
  C=(g_1,g_2,\ldots,g_m).
\]
单比特门只依赖其作用的一个逻辑量子比特；双比特门依赖两个逻辑量子比特的当前状态。
如果两个门作用在同一逻辑量子比特上，后一个门必须等待前一个门完成。这些先后关系
构成线路依赖结构。路由算法中的“前沿门”正是依赖已经满足、下一步可以尝试执行的门。

线路结构影响路由难度。GHZ 线路通常形成链式或星形纠缠，QFT 线路包含密集的长距离
双比特交互，QAOA 和 VQE 线路则由问题图或变分层决定交互模式。若交互图和硬件拓扑
形状接近，路由器可以少插入 SWAP；若逻辑交互跨越拓扑远端，编译器必须频繁移动
量子态。后续实验选择这些线路类型，就是为了覆盖不同交互结构下的路由行为。

\subsubsection*{硬件拓扑}

真实芯片的物理连接用图表示：
\[
  G=(V,E),
\]
其中 $V$ 是物理量子比特集合，$E$ 是可直接执行双比特门的耦合边。若
$(u,v)\in E$，映射到 $u$ 和 $v$ 上的两个逻辑量子比特可以直接执行双比特门；若
$(u,v)\notin E$，编译器必须先通过 SWAP 改变映射。

图 \ref{fig:hardware-topology} 给出本报告使用的两类拓扑示意。20 比特代表实验使用
Tokyo 风格耦合图，扩展实验使用二维网格。两类拓扑的共同点是连接稀疏：每个物理
量子比特只和少数邻居相连，远距离逻辑交互必须通过多次交换完成。系统会预先计算
物理距离 $d_G(u,v)$，即两个物理量子比特在耦合图中的最短路径长度，并把它作为
候选 SWAP 评分、前沿门代价和初始映射选择的重要依据：
\[
  d_G(u,v)=\mathrm{dist}_G(u,v).
\]

\begin{figure}[H]
  \centering
  \resizebox{\linewidth}{!}{%
  \begin{tikzpicture}[
  q/.style={circle, draw=blue!55!black, fill=blue!6, minimum size=0.58cm, font=\scriptsize},
  gridq/.style={circle, draw=gray!65!black, fill=gray!7, minimum size=0.48cm, font=\tiny},
  edge/.style={draw=gray!65, thick},
  label/.style={font=\small, align=center},
  note/.style={rectangle, draw=blue!40!black, fill=blue!3, rounded corners=2pt,
    text width=3.15cm, align=left, inner sep=5pt, font=\scriptsize}
]
  \node[label] at (1.2,2.35) {20Q 拓扑};
  \foreach \i/\x/\name in {0/0/v_0,1/1.1/v_1,2/2.2/v_2}
    \node[q] (a\i) at (\x,1.6) {$\name$};
  \foreach \i/\x/\name in {0/0/v_3,1/1.1/v_4,2/2.2/v_5}
    \node[q] (b\i) at (\x,0.65) {$\name$};
  \foreach \i/\x/\name in {0/0/v_6,1/1.1/v_7,2/2.2/v_8}
    \node[q] (c\i) at (\x,-0.3) {$\name$};
  \foreach \u/\v in {a0/a1,a1/a2,b0/b1,b1/b2,c0/c1,c1/c2,
    a0/b0,a1/b1,a2/b2,b0/c0,b1/c1,b2/c2,a0/b1,b1/c2}
    \draw[edge] (\u) -- (\v);

  \node[label] at (6.1,2.35) {网格拓扑};
  \foreach \i in {0,...,4} {
    \foreach \j in {0,...,2} {
      \node[gridq] (g\i\j) at (4.5+0.65*\i,0.35+0.65*\j) {};
    }
  }
  \foreach \i in {0,...,3} {
    \foreach \j in {0,...,2} {
      \pgfmathtruncatemacro{\next}{\i+1}
      \draw[edge] (g\i\j) -- (g\next\j);
    }
  }
  \foreach \i in {0,...,4} {
    \foreach \j in {0,...,1} {
      \pgfmathtruncatemacro{\next}{\j+1}
      \draw[edge] (g\i\j) -- (g\i\next);
    }
  }

  \node[note] at (9.35,1.2) {
    顶点表示物理量子比特，边表示可直接执行双比特门的耦合关系。路由器只允许在边上
    执行目标门或插入 SWAP。
  };
\end{tikzpicture}

  }
  \caption{硬件拓扑}
  \label{fig:hardware-topology}
\end{figure}

\subsection{映射与路由目标}

逻辑量子比特集合记为 $Q$。路由过程中的映射写作
\[
  \pi:Q\rightarrow V ,
\]
表示逻辑量子比特当前被放置到哪个物理位置。若双比特门
$g=(q_i,q_j)$ 满足
\[
  (\pi(q_i),\pi(q_j))\in E,
\]
该门可以在当前映射下直接执行。若条件不满足，路由器需要选择一条物理边 $(u,v)$
插入 SWAP，并交换这两个物理位置上承载的逻辑量子比特。设 $\pi'$ 是执行 SWAP 后的
新映射，则位于 $u$ 和 $v$ 的两个逻辑量子比特互换位置，其余逻辑量子比特保持不变。

图 \ref{fig:routing-model} 用一个小例子说明映射和 SWAP 的关系。逻辑线路中
$q_0$ 和 $q_2$ 需要执行 CNOT，但初始映射把它们放在不相邻的物理位置。编译器不能
直接执行该门，只能先在硬件边上插入 SWAP，让逻辑量子态沿拓扑移动。这个例子解释了
为什么路由不是简单地把每个逻辑门翻译成一个物理门：一次 SWAP 会改变后续所有门看到
的物理位置，当前选择和未来代价紧密耦合。

\begin{figure}[H]
  \centering
  \begin{tikzpicture}[
  font=\small,
  title/.style={font=\bfseries\small, align=center},
  note/.style={font=\scriptsize, align=center, text=black!72},
  qnode/.style={circle, draw=blue!60!black, fill=blue!7, minimum size=0.56cm, font=\scriptsize},
  vnode/.style={circle, draw=black!55, fill=gray!6, minimum size=0.58cm, font=\scriptsize},
  occ/.style={circle, draw=green!45!black, fill=green!10, minimum size=0.58cm, font=\scriptsize},
  gate/.style={rectangle, draw=blue!55!black, fill=blue!5, rounded corners=2pt,
    minimum width=1.34cm, minimum height=0.52cm, align=center, font=\scriptsize},
  swap/.style={rectangle, draw=orange!70!black, fill=orange!10, rounded corners=2pt,
    minimum width=1.34cm, minimum height=0.52cm, align=center, font=\scriptsize},
  arrow/.style={-{Latex[length=2mm]}, thick, draw=black!62},
  edge/.style={thick, draw=black!42}
]
  \node[title] at (-3.95,2.05) {逻辑线路};
  \node[qnode] (q0) at (-5.25,1.15) {$q_0$};
  \node[qnode] (q1) at (-4.15,1.15) {$q_1$};
  \node[qnode] (q2) at (-3.05,1.15) {$q_2$};
  \node[gate] (g1) at (-4.15,0.25) {目标门\\$CX(q_0,q_2)$};
  \draw[edge, blue!45!black] (q0) -- (q1);
  \draw[edge, blue!45!black] (q1) -- (q2);
  \draw[edge, blue!70!black] (q0) to[bend left=32] (q2);

  \node[title] at (0,2.05) {初始映射};
  \node[occ] (a0) at (-1.15,1.18) {$v_0/q_0$};
  \node[vnode] (a1) at (0,1.18) {$v_1$};
  \node[occ] (a2) at (1.15,1.18) {$v_2/q_2$};
  \draw[edge] (a0) -- (a1);
  \draw[edge] (a1) -- (a2);
  \node[note] at (0,0.27) {$q_0$ 与 $q_2$ 不相邻，\\双比特门不能直接执行};

  \node[title] at (3.95,2.05) {插入 SWAP};
  \node[vnode] (b0) at (2.80,1.18) {$v_0$};
  \node[occ] (b1) at (3.95,1.18) {$v_1/q_0$};
  \node[occ] (b2) at (5.10,1.18) {$v_2/q_2$};
  \draw[edge] (b0) -- (b1);
  \draw[edge] (b1) -- (b2);
  \draw[edge, orange!80!black, line width=1.2pt] (b0) -- node[below=3pt, font=\scriptsize] {SWAP} (b1);
  \node[note] at (3.95,0.27) {交换后 $q_0$ 与 $q_2$ 相邻，\\原始门可在硬件边上执行};

  \draw[arrow] (-2.35,1.16) -- (-1.55,1.16);
  \draw[arrow] (1.55,1.16) -- (2.35,1.16);
  \node[swap] at (0,-0.86) {路由目标：在保持原线路语义的前提下，减少额外 SWAP 和线路深度};
\end{tikzpicture}

  \caption{路由问题}
  \label{fig:routing-model}
\end{figure}

\subsubsection*{评价指标}

本项目使用三个主要指标衡量路由结果：SWAP 数、线路深度和编译耗时。SWAP 数表示
为满足拓扑约束额外插入了多少交换门，它直接反映路由带来的额外操作。线路深度表示
按依赖关系分层后需要多少时间步执行，深度越大，量子态在噪声环境中停留的时间越长。
编译耗时表示算法生成结果需要的时间，它决定在线演示、服务器接口和本地工具的响应
速度。

三者可以写成一个综合优化目标：
\[
  \min \; \alpha N_{\mathrm{swap}}+\beta D+\gamma T ,
\]
其中 $N_{\mathrm{swap}}$ 是插入 SWAP 数，$D$ 是路由后线路深度，$T$ 是编译耗时，
$\alpha,\beta,\gamma$ 用于表示不同指标的权重。实际系统不求解精确全局最优解，而是
在有限时间内寻找拓扑合法、语义等价、质量较好的可行线路。这个选择符合量子信息基础大作业的
实验目标：重点比较不同算法设计在相同线路和拓扑下的可执行结果，而不是把问题化为
一个难以在报告规模内完整求解的精确优化模型。

\subsection{课程知识对应}

表 \ref{tab:concept-mapping} 概括了课程知识和本项目实现之间的对应关系。该表的作用
不是把课程概念逐条罗列，而是说明路由编译为什么适合作为量子信息基础大作业主题：
它同时涉及量子态、酉门、纠缠线路、硬件约束、噪声影响和近似优化。

\begin{table}[H]
  \centering
  \caption{课程对应}
  \label{tab:concept-mapping}
  \begin{tabularx}{\linewidth}{lY}
    \toprule
    知识点 & 项目体现 \\
    \midrule
    量子态 & 线路中的门序列描述量子态演化，路由输出必须保持逻辑语义等价。\\
    张量积 & 多量子比特寄存器位于张量积空间，双比特门作用在两个子系统上。\\
    酉门 & CNOT、CZ、SWAP 都是酉门；额外 SWAP 改变承载位置，不改变理想演化。\\
    纠缠 & GHZ、QFT 等样例包含密集双比特交互，对硬件拓扑更敏感。\\
    拓扑 & 芯片耦合关系被建模为图，只有图边上的物理量子比特能直接执行双比特门。\\
    噪声 & 额外 SWAP 和更深线路会增加噪声暴露时间，减少交换门具有物理意义。\\
    优化 & 映射选择、候选 SWAP、搜索剪枝构成受约束的组合优化过程。\\
    验证 & 轨迹复放逐步检查门操作和映射变化，确认输出线路满足拓扑约束。\\
    \bottomrule
  \end{tabularx}
\end{table}

